\section{Flags and decompositions in a cyclic family}
\label{sec:companion-family}

A proper flag and a direct-sum decomposition agree on a cyclic module
with distinct support, but they have different extensions across a
collision. We compute this difference for the length-three family.
The calculation concerns the actual ordinary direct images and their
base-change map. It also explains why the proper flag trace cannot be
recovered from the ordinary decomposition image.

Throughout this section the coefficient is the normalized constant
rational coefficient of \eqref{gh:relative-coefficient}. Fix a coordinate
disk $D$ in one common elliptic torsor and put
\[
 X=D^3,\qquad Y=\operatorname{Sym}^3D,\qquad m:X\longrightarrow Y.
\]
Let $U\subset X$ be the ordered distinct-root locus, let $V=m(U)$,
and write $j:U\hookrightarrow X$ and $j_V:V\hookrightarrow Y$.
The finite map $m$ has degree six. On $Y$ the universal cyclic module is
\[
 \mathcal T=\mathcal O_Y[\zeta]/(f_{\mathrm{univ}}(\zeta)).
\]
Its pullback to $X$ is
$\mathcal O_X[\zeta]/\prod_{i=1}^3(\zeta-t_i)$.
The variables $t_i$ are support coordinates; they are distinct from
the stabilizer Chern classes $z_i$.

Define ordinary derived direct images and their restriction map by
\begin{equation}\label{cf:images}
 \mathcal F=Rm_*\mathbb Q_X,\qquad
 \mathcal G=R(mj)_*\mathbb Q_U,\qquad
 r:\mathcal F\longrightarrow\mathcal G.
\end{equation}
The map $r$ is induced by $\mathbb Q_X\to Rj_*\mathbb Q_U$.
To distinguish a labelled family from the ordered cover defining
these images, let $T$ be a second copy of $X$ and write $g=m:T\to Y$.
Set
\[
 Z=T\times_YX,\quad p:Z\to T,\quad
 \mathcal D_T=T\times_YU,\quad
 \mathcal F_T=Rp_*\mathbb Q_Z,\quad
 \mathcal G_T=R(\mathcal D_T\to T)_*\mathbb Q.
\]
Proper base change identifies $\mathcal F_T$ with $g^*\mathcal F$.
There is an ordinary base-change transformation
\begin{equation}\label{cf:base-change}
 b:g^*\mathcal G\longrightarrow\mathcal G_T.
\end{equation}
The construction and adjunctions for these derived images are the
usual ones for rational constructible complexes
\cite[\S\S3.4,4.1.4,4.3.2--4.3.3]{deCataldoMigliorini2007}.

\subsection{The two correspondence families}

\begin{proposition}\label{cf:families}
The complete proper flag family of $\mathcal T$ is $m:X\to Y$,
including its nonreduced fibres. Its complete ordered direct-sum
decomposition family is $mj:U\to Y$.
The two parenthesizations of either complete correspondence are
canonically identified. After pullback to $T$, the decomposition
family is
\[
 \mathcal D_T\simeq\coprod_{\sigma\in S_3}U,
\]
where the component $\sigma$ assigns ordered endpoints
$t_{\sigma(1)},t_{\sigma(2)},t_{\sigma(3)}$.
\end{proposition}
\begin{proof}
For a cyclic module $A=R[\zeta]/(f)$, every endomorphism is
multiplication by an element of $A$. A decomposition therefore gives
orthogonal idempotents. If the three roots are distinct, the Chinese
remainder theorem gives the three line summands in any prescribed
order. At a repeated root the corresponding local algebra has no
nontrivial idempotent. Its cyclic Jordan block cannot be a sum of
length-one modules at that point.

This argument also identifies the family, rather than only its closed
points. Project the cyclic generator $1$ onto the three rank-one
endpoints. It generates each endpoint and hence trivializes it.
In these trivializations the images of $1,\zeta,\zeta^2$ form a
Vandermonde matrix. It is invertible precisely when every difference
of two endpoint roots is a unit. Chinese remainders then give the
inverse decomposition on exactly that open locus. Thus the Cartesian
pullback of the decomposition correspondence, with its identification
of the middle object with $\mathcal T$, is $U$. Endpoint automorphisms
are cancelled by that identification; no additional
$B\mathbb G_m$ factors remain.

Flags require no idempotents. A monic factorization $f=gh$ gives
\[
 0\longrightarrow (g)/(gh)\longrightarrow R[\zeta]/(gh)
       \longrightarrow R[\zeta]/(g)\longrightarrow0.
\]
Conversely, a locally free quotient of a cyclic module is cyclic.
Its characteristic polynomial is monic; division by this polynomial
identifies its kernel with the corresponding ideal. Applying this
argument successively identifies a full flag with an ordered monic
factorization into three linear factors. These constructions commute
with base change and include infinitesimal factorizations. Their
parameter scheme is $X$ over $Y$. In particular, the flag fibre at
$3[0]$ has one geometric point and scheme length six.

Both iterated decompositions specify the same three ordered orthogonal
idempotents. Both proper flag refinements specify the same ideal
filtration. The resulting identifications hold on the schemes of
families. Composition of ordinary direct images identifies the two
decomposition images; composition of proper traces identifies the two
flag traces. Finally, pulling the ordered-root covering back to $T$
over $U$ gives its six permutation graphs, as asserted.
\end{proof}

\subsection{Ordinary stalks and ramified base change}

Write $y_0=3[0]\in Y$ and $t_0=(0,0,0)\in T$. All cohomological
degrees in the following statement are ordinary cochain degrees.

\begin{theorem}\label{cf:comparison}
The central stalk of $\mathcal F$ is $\mathbb Q$ in degree zero.
The central stalk of $\mathcal G$ has cohomology algebra
\begin{equation}\label{cf:arrangement-ring}
 H=\Lambda_{\mathbb Q}(a,b,c)/(ab-ac+bc),\qquad
 |a|=|b|=|c|=1.
\end{equation}
Its dimensions in degrees $0,1,2$ are $(1,3,2)$, and $r$ sends
the degree-zero unit to the unit. The central cohomology dimensions
of $\mathcal G_T$ are $(6,18,12)$. Under these identifications,
\begin{equation}\label{cf:actual-base-change}
 H^*(b_{t_0}):H\longrightarrow\bigoplus_{\sigma\in S_3}H,
 \qquad v\longmapsto(\sigma^*v)_{\sigma\in S_3},
\end{equation}
where $S_3$ permutes the three unordered edge classes $a,b,c$.
Consequently the cohomology dimensions of $\operatorname{Cone}(b)$
at $t_0$ are $(5,15,10)$ in degrees $0,1,2$. Those of the cone of
\begin{equation}\label{cf:composite}
 \mathcal F_T\simeq g^*\mathcal F
       \xrightarrow{g^*r}g^*\mathcal G
       \xrightarrow{b}\mathcal G_T
\end{equation}
are $(5,18,12)$ in those degrees. Neither map is an equivalence.
\end{theorem}
\begin{proof}
The inverse images
$m^{-1}(\operatorname{Sym}^3D_\epsilon)=D_\epsilon^3$ form
cofinal neighbourhoods of $m^{-1}(y_0)$. Intersecting with $U$
gives a central hyperplane-arrangement complement in a polydisk.
Positive radial rescaling identifies its homotopy type with that
of the full complement. Thus the neighbourhood definition of an
ordinary direct-image stalk gives
\[
 \mathcal G_{y_0}\simeq
 R\Gamma(\operatorname{Conf}_3(\mathbb C),\mathbb Q).
\]
No proper base-change statement has been applied to $mj$.

The coordinate change
\[
 (t_1,t_2,t_3)\longmapsto
 \left(t_1,\ x=t_2-t_1,\ v=\frac{t_3-t_1}{t_2-t_1}\right)
\]
identifies this configuration space with
$\mathbb C\times\mathbb C^*\times(\mathbb C\setminus\{0,1\})$.
The last factor retracts onto a wedge of two circles. Let $A$ be
the degree-one generator of $\mathbb C^*$ and let $P,Q$ be the
two punctured-line generators. Normalize logarithmic classes by
$1/(2\pi i)$ and set
\[
 a=[d\log(t_2-t_1)],\quad
 b=[d\log(t_3-t_1)],\quad
 c=[d\log(t_3-t_2)].
\]
Then $a=A$, $b=A+P$, and $c=A+Q$. The relations
$A^2=P^2=Q^2=PQ=0$ give $ab-ac+bc=0$.
These classes generate the product cohomology, whose dimensions
are $(1,3,2)$, proving \eqref{cf:arrangement-ring}.
Reversing an edge multiplies its difference by the constant $-1$
and hence leaves its logarithmic class unchanged. The permutation
action therefore permutes the three unordered edges and induces
the two-dimensional standard representation in degree two.

Proper base change applies to $m$. Its central fibre has one
underlying point, so $\mathcal F_{y_0}=\mathbb Q$. Restriction
preserves the unit. In particular, the cone of $r$ has dimensions
three and two in degrees one and two.

Proposition~\ref{cf:families} gives
$\mathcal G_T=\bigoplus_{\sigma\in S_3}Rj_*\mathbb Q_U$.
On every small configuration neighbourhood the component of the
base-change map labelled by $\sigma$ is pullback by that root
permutation. This proves \eqref{cf:actual-base-change}; replacing
$\sigma$ by its inverse merely reorders the summands.
The identity component makes this map injective in every degree,
with ranks $(1,3,2)$. Its cokernel dimensions are therefore
$(5,15,10)$. The source of \eqref{cf:composite} has only its
degree-zero unit, which maps to the diagonal unit in six summands.
Its cone has dimensions $(5,18,12)$.
\end{proof}

The unlabelled nearby covering is connected, while its labelled
strict pullback has six components. Neither the empty strict central
decomposition fibre nor a six-dimensional degree-zero vector space
computes these ordinary derived images.

\subsection{The proper collision coefficient and its trace}

Let $H_{ij}\subset T$ be the diagonal $t_i=t_j$.
The reduced space underlying $Z$ is the union of its six permutation
graphs. On $H_{ij}$, the labels $\sigma$ and $(ij)\sigma$ have the
same root values. Define a sheaf on $T$ by
\begin{equation}\label{cf:collision-kernel}
 \mathcal K=\ker\left[
 \bigoplus_{\sigma\in S_3}\mathbb Q_T
 \xrightarrow{\ d\ }
 \bigoplus_{i<j}\ \bigoplus_{\{\sigma,(ij)\sigma\}}
                         \mathbb Q_{H_{ij}}\right].
\end{equation}
Each diagonal contributes three unordered pairs of labels.
Here $\mathbb Q_{H_{ij}}$ denotes closed direct image to $T$,
and each component of $d$ is the difference of the two entries
restricted to that diagonal. No surjectivity of $d$ is asserted.

\begin{proposition}\label{cf:proper-coefficient}
There is a canonical isomorphism $\mathcal F_T\simeq\mathcal K$
in degree zero. Its ranks are six at distinct roots, three at a
generic double collision, and one at the triple collision.
The proper trace and unit are
\begin{equation}\label{cf:trace}
 \mathcal K\hookrightarrow\bigoplus_{\sigma\in S_3}\mathbb Q_T
       \xrightarrow{\ \sum\ }\mathbb Q_T,
 \qquad
 \mathbb Q_T\longrightarrow\mathcal K,\quad
       1\longmapsto(1)_\sigma.
\end{equation}
Their composite is $6\operatorname{id}_{\mathbb Q_T}$.
Both proper parenthesizations retain these maps and all their
collision restrictions.
\end{proposition}
\begin{proof}
Pull a locally constant function on $Z$ to its six graph components.
The entries agree on every graph intersection. These agreements
are generated by the transpositions in \eqref{cf:collision-kernel}.
Conversely, compatible local constants glue on the finite union
of closed graphs. This identifies $p_*\mathbb Q_Z$ with
$\mathcal K$. The finite proper fibres have no higher rational
cohomology, so $Rp_*\mathbb Q_Z=p_*\mathbb Q_Z$.

At a double collision the labels form three pairs, and at the triple
collision they all agree. The sum in \eqref{cf:trace} thus weights
the three values by two at a double collision and the single value
by six at the triple collision. These are exactly the pullbacks of
the finite proper trace: its local degree is the multiplicity of
the corresponding sheets. The diagonal unit proves the composite
formula. Proposition~\ref{cf:families} identifies both refinements
as the same proper map, so functoriality of the trace gives the
last assertion.
\end{proof}

The partial ordered decomposition of type $1+2$ isolates one root.
Its three components on $T$ are
\[
 U_i=\{t_i\ne t_j\text{ for every }j\ne i\},\qquad i=1,2,3.
\]
Writing $j_i:U_i\hookrightarrow T$, its ordinary image is
$\bigoplus_iR(j_i)_*\mathbb Q_{U_i}$.
Each central local complement has homotopy type $(\mathbb C^*)^2$,
so the central cohomology dimensions are $(3,6,3)$.
The proper partial flag has a one-dimensional central stalk and
generic degree three; refining its length-two factor has degree two.
Both proper refinements therefore give the same trace of degree six
and the same collision agreements in \eqref{cf:collision-kernel}.

The endpoint lines retain their quotient-first order. For an
invertible cyclic coefficient $N$, let $e_i$ be evaluation at the
$i$th ordered root and put $\ell_i=c_1(e_i^*N)$. The lines are
\begin{equation}\label{cf:endpoint-lines}
 L_i=e_i^*N\otimes
       \mathcal O\left(-\sum_{j<i}\Delta_{ji}\right).
\end{equation}
Indeed, each successive ideal quotient contributes the earlier
linear factors; restriction to the next graph gives the displayed
diagonal divisors. Thus a polynomial insertion $P$ has cyclic trace
\[
 m_*P(\ell_1,\ell_2-\Delta_{12},
                \ell_3-\Delta_{13}-\Delta_{23}).
\]
In particular, the third endpoint contains both preceding diagonals.
Although the local support diagonals are principal and their ordinary
divisor Chern classes can vanish on a polydisk, their defining sections,
the incidence maps of \eqref{cf:collision-kernel}, the logarithmic
classes of \eqref{cf:arrangement-ring}, and the global endpoint
formula remain distinct pieces of the geometry.

\subsection{The obstruction to extending the trace}

\begin{theorem}\label{cf:no-trace}
There is no degree-zero morphism $\mathcal G\to\mathbb Q_Y$
extending the generic proper flag trace. There is likewise no
morphism $\mathcal G_T\to\mathbb Q_T$ extending generic summation
over the six ordered decompositions. Consequently the proper traces
cannot factor through $r$ or through \eqref{cf:composite}.
\end{theorem}
\begin{proof}
On $V$ let $\mathcal L=(m|_U)_*\mathbb Q_U$, the rank-six
permutation local system. Then
$\mathcal G=R(j_V)_*\mathcal L$.
The unit $\eta:\mathbb Q_Y\to\mathcal G$ restricts on $V$ to
the vector of six ones. If $t:\mathcal G\to\mathbb Q_Y$
extended summation, $t\eta$ would restrict to $6\operatorname{id}$.
Since $Y$ is connected and
$\operatorname{End}(\mathbb Q_Y)=H^0(Y;\mathbb Q)=\mathbb Q$,
we would have $t\eta=6\operatorname{id}_{\mathbb Q_Y}$.

Choose a closed point $i:\{y\}\hookrightarrow Y\setminus V$.
For every complex $C$ on the point, adjunction gives
\[
 \begin{split}
 \operatorname{Hom}(C,i^!R(j_V)_*\mathcal L)
 &\simeq\operatorname{Hom}(i_*C,R(j_V)_*\mathcal L)\\
 &\simeq\operatorname{Hom}(j_V^*i_*C,\mathcal L)=0.
 \end{split}
\]
Hence $i^!\mathcal G=0$. On the other hand, $Y$ is a smooth
complex threefold, and relative cohomology of a small real six-ball
and its punctured ball gives $i^!\mathbb Q_Y=\mathbb Q[-6]$.
Applying $i^!$ to $t\eta$ would factor the nonzero map
$6\operatorname{id}_{\mathbb Q[-6]}$ through zero, a contradiction.
The same argument on $T$ uses
$\mathcal G_T=Rj_*(\mathbb Q_U^{\oplus6})$ and proves the
labelled assertion.
\end{proof}

This obstruction is a statement about morphisms of constructible
complexes, not only stalk dimensions. A linear projection from the
central cohomology $(6,18,12)$ to $\mathbb Q$ does not supply the
missing trace. The proper coefficient $\mathcal F$, and its explicit
pullback $\mathcal K$, retain the trace but are not equivalent to
the ordinary decomposition images.

Theorem~\ref{cf:no-trace} concerns the strict pullback fixing a
cyclic middle object. The next section constructs the coefficient
sheaf over support and its direct-sum coalgebra.
Section~\ref{sec:supported-springer} constructs the proper
Springer action and the supported Hall Hopf comparison, then
compares their evaluations with $\mathcal F$ and $\mathcal G$.
These maps retain the nonsplit cyclic fibre and do not extend
the forbidden trace.
